% 第 7 章给作者的建议（section 由主控文件给出）

本章把前几章的推导与核验结果转化为可操作的修改方案、答辩口径与后续研究计划。总体判断如下。六个命题与推论 1 的代数结论全部成立，第~\ref{sec:props} 节逐式补齐的推导没有发现错误。第~\ref{sec:checks} 节汇总的 199 项核验中，187 项通过、10 项存在偏差、2 项未通过，两项未通过分别源于一次蒙特卡洛抽样波动和原文第 508 行的措辞错误。需要修改之处集中在三类问题：近似的方向与误差没有交代，若干带有严格、下界、相符等字眼的表述过强，实证截距与斜率到理论参数的映射缺少中间环节。风险最高的四处分别是截距 $a$ 的口径、斜率 $b$ 的随机分母放大、第 508 行的“相符”以及第 393、406 行的“下界”说法。其中随机分母机制可以把滞后权重敏感性从需要辩解的弱点转为理论能够定量预测的结果。

\subsection{论文修改建议}\label{sub:ad-revise}

下文每条建议依次给出原文位置、存在的问题、建议改法与依据。行号均指论文源码的行号，依据指向本讲义的相应章节。三组建议分别对应严谨性、实证与写作图表，严谨性一组按优先级排序。

\subsubsection{严谨性修改（按优先级排序）}

\begin{enumerate}
\item 改正三处不成立的表述。这三处都会在答辩或审稿中被直接追问，修改只涉及措辞，不需要新的计算。
\begin{itemize}
\item 原文第 508 行称截距“经验值 0.0991 与参数空间隐含值 0.014--0.029 相符”，第 515 行称比值“数量级上与之相符”。按第 515 行的流感情景计算，实证 $a$ 是情景隐含值的 3.5--6.9 倍，$b=2.2434$ 是隐含值 0.610--0.705 的 3.2--3.7 倍；$b/a=22.6$ 落入理论区间，是因为两者被放大了相近的倍数。建议第 508 行删去“相符”，改为“高于情景隐含值（3.5--6.9 倍），差额来自周增长比与世代再生数的口径差及代理偏差”；第 515 行改为“$b/a$ 与情景区间同量级，但 $a$、$b$ 各自均高于隐含值，这一吻合不构成结构识别”。依据见第~\ref{sec:empirical} 节的情景隐含值计算。
\item 原文第 393、406 行称按实际暴露规模考察时，经验预测误差“始终不低于单代下界”。条件偏差--方差分解（原文式 28）只保证偏差项非负，推不出总误差不低于 $v_R+A/m$；按初始规模排列时，经验误差与参照之比为 0.75--0.98，按实现规模分箱后才为 1.01--1.17。建议改为“按实现规模分箱时，本仿真中经验误差不低于参照曲线（1.01--1.17 倍），方向与非负偏差项一致，但这并非一般性下界”。依据见第~\ref{sec:props} 节对更新过程分解的讨论。
\item 原文第 512 行称分箱 $\mathrm{CV}^2$ 的递减“为 FGLS 采用 $1/\hat\mu^2$ 权重提供了依据”，论证方向相反。$1/\hat\mu^2$ 成为最优权重的条件是 $\mathrm{CV}^2$ 为常数，而分箱 $\mathrm{CV}^2$ 从 5.32 降到 0.99，恰好违背这一条件。建议改写为“$\mathrm{CV}^2$ 随规模下降，表明 $1/\hat\mu^2$ 并非最优权重；在均值设定正确时，任意只依赖 $D$ 的权重均给出一致估计，权重只影响效率”。依据见第~\ref{sec:empirical} 节的 FGLS 一致性论证与第~\ref{sec:checks} 节的仿真核验。
\end{itemize}

\item 在实证节加入随机分母的定量说明。原文第 506--512 行把斜率 $b$ 直接对应为 $A/\bar R^2$，没有考虑分母 $D_t$ 本身由报告病例构造、带有二项抽样噪声。由一阶 delta 法，斜率的近似表达式如下：
\begin{equation}\label{eq:ad-denom}
b\approx\frac{A+(1-\rho)(\bar R^2+v_R)\sum_s w_s^2}{\bar R^2}.
\end{equation}
式~(\ref{eq:ad-denom}) 同时解释了三个现象。第一，三种滞后权重的 $\sum_s w_s^2$ 分别为 0.495、1 与 1/3，预测的 $b$ 比值为 1.44 与 0.86，实证比值为 1.40 与 0.86。第二，在流感情景参数下，忽略分母噪声时理论 $b$ 只有实证值的 29\%，计入二项分母噪声后升至 51\%；若报告方差存在约 $\varphi\approx1.96$ 倍的过离散，剩余缺口即可补足，这是一种可能的解释，尚需数据检验。第三，分母噪声只抬高斜率、不改变截距，所以 $b/a$ 会高估 $m_\times$，模拟中 $b/a$ 为 59.2，而对应的 $A/v_R$ 为 30.7。建议把这一机制写入局限性一节，并作为新增结果正面报告，替代现有“经验拐点对回溯窗口敏感”的单纯描述。依据见第~\ref{sec:empirical} 节随机分母的推导与第~\ref{sec:checks} 节的蒙特卡洛核对。

\item 改写截距 $a$ 的解释。原文第 441--446 行构造的周增长比 $G_{j,t+1}=C_{j,t+1}/D_{j,t}$ 跨越两个以上的流感世代（世代间隔约 3 天，平均回溯 1.45 周），理论中的 $R_t$ 却是单代量。若 $a$ 只是环境方差 $v_R/\bar R^2$，它应与滞后权重无关；实证 $a$ 却随平均滞后单调增大，单周滞后（1.0 周）、原始权重（1.45 周）与三周等权（2.0 周）下分别为 0.0682、0.0991 与 0.1524。建议在第 446 行后说明 $G$ 与 $R_t$ 的口径差，把 $a$ 解读为“周增长比的规模无关离散”，并把三种权重下的 $a$ 列为证据。依据见第~\ref{sec:empirical} 节。

\item 集中写出模型假设。原文第 91--94 行只给出 $X_{i,t}$ 的前两阶矩，附录第 688 行却直接使用完整的负二项分布；命题 1 的精确分布、极大似然估计与 CRLB，以及命题 6 的精确检验都依赖完整分布。建议在第 94 行后增设假设 1，列出四条内容：个体继发感染给定 $R_t$ 独立同分布于 $\mathrm{NB}(k,k/(k+R_t))$，等价于 Poisson--Gamma 混合；同期 $n$ 名感染者共享同一 $R_t$；$k$ 与 $\rho$ 已知；二项报告与传播条件独立。第 114 行之后再补一句，说明 $\mathbb{E}[R_t\mid n]=\bar R$ 对设计固定的 $n$ 成立，对由过去传播决定的 $n_t$ 需另行假设，因为环境存在自相关时，较大的 $n_t$ 意味着 $R_{t-1}$ 偏大，进而 $R_t$ 也偏大。第~\ref{sec:model} 节已整理这些假设，附录引理~\ref{lem:ap-nb} 说明了共享 $R_t$ 在负二项可加性中的作用。

\item 修正命题 1 的表述并补齐 UMVUE 论证。原文第 145 行、第 196 行图注与第 419 行把 CRLB 写成有限样本下的“收敛”或“紧致性”。得分函数恰好等于 $\mathcal I(R)(\widehat R-R)$，所以等号对每个 $n$ 都成立；图 1(b) 中 0.988--1.017 的比值是 50,000 次重复下的蒙特卡洛误差。建议第 419 行改为“仿真方差与理论值之比处于 0.988--1.017，偏离在蒙特卡洛误差范围内，与 CRLB 精确达到一致”。附录第 715 行应补一句 Lehmann--Scheffé 论证：固定 $nk$ 时负二项族是关于 $C_t$ 的单参数指数族，$C_t$ 为完备充分统计量，$C_t/(\rho n)$ 无偏，故为 UMVUE \citep{lehmann1998point}。两条论证都要求 $k$ 已知，第~\ref{sec:props} 节给出逐步推导，附录引理~\ref{lem:ap-fisher} 给出取等条件的完整证明。

\item 修正命题 4 与推论 1 的近似与数值。第 254--256 行的 $T_{\rm eff}$ 近似恒高估环境项，相对高估为 $2r(1-r^T)/[T(1-r)(1+r)-2r(1-r^T)]$，$T=26$、$r=0.8$ 时为 20.5\%，$T=10$、$r=0.8$ 时为 65.8\%，所以条件 $T\gg1/(1-r)$ 不足以保证近似良好。建议正文直接给出附录第 751 行的精确闭式，$T_{\rm eff}$ 只作直觉保留；或把条件改为 $T\gg2r/(1-r^2)$，并写明近似偏保守。第 283 行的 14.6--22.2 与 46.1--70.2 来自第 515 行流感情景的 $m_\times$ 区间，正文应写明来源；情景上端按一致参数为 49.18，49.3 应改为 49.2，第 515 行同步修改。另需补充 $T$ 为整数的说明：$c_0=100$、$c_1=1$、$\mathcal B=600$ 时连续解为 $m^*=72.86$，整数最优为 $T=3$、$m=100$。AR(1) 修正同样受影响，精确整数搜索给出 42.9（$r=0.5$）与 25.0（$r=0.8$），近似公式给出 42.1 与 24.3。依据见第~\ref{sec:props} 节、图~\ref{fig:pb-ar1} 与图~\ref{fig:pb-budget}。

\item 修正命题 6 的检验表述。第 340 行“实际检验水平严格低于名义水平 $\alpha$”应改为“不超过”，核验网格内实际水平最大达到 0.05000。非随机化检验只在它自己的实际水平下是一致最优；若需要恰好水平 $\alpha$，应在命题后注明可用随机化检验 \citep{karlin1956theory,lehmann2005testing}。局限性中应讨论以 $\hat k$ 代入后检验水平失去精确保证。图 2（第 358 行）报告的是功效首次达到 0.8 的 $n$，即 2071、568、229、568、1660，此后功效会跌回 0.8 以下；此后恒不低于 0.8 的稳定样本量为 2107、581、235、585、1684，多出 1.4\%--3.0\%，建议同时报告。第 364 行的 Wald 检验水平可以精确计算，$m=5$、$k=0.1$ 时为 0.0781，论文的蒙特卡洛值为 8.1\%，建议改报精确值。依据见第~\ref{sec:props} 节与图~\ref{fig:pb-power}。

\item 为命题 5 补可计算的展开。原文第 299--308 行的两式是对任意分布成立的全方差恒等式，命题本身没有给出可计算的量。在负二项分支且 $\Lambda_{t+1}$ 与 $R_{t+1}$ 条件独立的附加假设下，归一化的预测方差可展开如下：
\begin{equation}\label{eq:ad-prop5}
\frac{\operatorname{Var}(C_{t+1}\mid\mathcal F_t^C)}{(\rho E_\Lambda)^2}=v_R+\frac{A}{m}+\mathrm{CV}^2(\Lambda_{t+1}\mid\mathcal F_t^C)\,(\bar R^2+v_R),
\end{equation}
其中 $E_\Lambda$ 为 $\Lambda_{t+1}$ 的条件均值，$m=\rho E_\Lambda$。式~(\ref{eq:ad-prop5}) 在 $\Lambda_{t+1}$ 已知时退化为命题 2 的 $v_R+A/m$，多出的 $\mathrm{CV}^2$ 项正是“单代静态公式不能直接外推”的数学含义。建议把它作为命题 5 的推论写入正文，并删除第 312 行与上下文没有衔接的“瓶颈属于目标 B”一句，或改写为该推论的解释。推导见第~\ref{sec:props} 节，其中两次用到附录引理~\ref{lem:ap-totvar} 的条件全方差公式。

\item 为命题 3 补边际比较并消除前后矛盾。命题 3 只比较了两个极限。同比例增加期望报告数 $m$ 时，沿 $n$ 与沿 $\rho$ 两条路径的边际降幅之比如下：
\begin{equation}\label{eq:ad-ratio}
\frac{A}{\bar R}=1+\frac{\rho(\bar R^2+v_R)}{k\bar R}.
\end{equation}
图 1 参数下这一比值为 1.85，流感住院情景（$\rho\le0.021$）下只有 1.02--1.18。建议在命题 3 后加注记给出式~(\ref{eq:ad-ratio})，并在第 467、613、658 行说明住院监测中两条杠杆近似等价。第 168 行“无法随横截面聚合消除”的说法与第 617 行的全国汇总结论矛盾，应改为“无法通过同一环境内扩大规模消除；跨相关性不完全的环境聚合可削减其 $(1-\bar c)$ 部分”。依据见第~\ref{sec:props} 节。
\end{enumerate}

\subsubsection{实证部分的补充}

\begin{enumerate}
\item 补充两向聚类。按州聚类处理的是同州跨周的相关，同周跨州的共同冲击会破坏簇间独立，而论文目前只做了单向的按周聚类敏感性分析。建议按 \citet{cameron2011robust} 的两向聚类方法重新给出 $b$ 与 $m_\times$ 的区间。51 个州簇处在聚类推断可用的下限附近，小簇数下的 bootstrap 修正可参考 \citet{cameron2008bootstrap} 与 \citet{cameron2015practitioner}。核验模拟中州聚类 bootstrap 对 $b$ 与 $m_\times$ 的覆盖率均为 0.944，未见明显欠覆盖，但模拟只包含州间独立的情形。依据见第~\ref{sec:empirical} 节与第~\ref{sec:checks} 节。

\item 补充 FGLS 权重对照。建议增加 $1/(\hat\mu^2\widehat{\mathrm{CV}}{}^2(D))$ 权重，其中 $\widehat{\mathrm{CV}}{}^2(D)$ 对 $\log D$ 作线性拟合，或者给出 OLS 对照 \citep{wooldridge2010econometric}。核验模拟显示，在 $\mathrm{CV}^2$ 随规模递减的设定下，FGLS 的 $\hat a$ 与 $\hat b$ 无偏（$\hat b$ 低 2.4\%，在 1.2--1.7 个标准误之内），$\hat b/\hat a$ 也没有比值正偏；更有效的权重主要使区间收窄，点估计应基本不变。依据见第~\ref{sec:empirical} 节。

\item 报告过离散参数。用分箱的 $\operatorname{Var}(C)/\mathbb{E}[C]$ 或周内日度数据估计报告过离散参数 $\varphi$，检验 $\varphi\approx2$ 是否合理。若成立，计入过离散的式~(\ref{eq:ad-denom}) 推广形式即可闭合斜率水平的缺口；若不成立，需要寻找其他放大来源。依据见第~\ref{sec:empirical} 节。

\item 报告空间汇聚的隐含相关。第 617 行的全国汇总只有 $N=28$ 个观测，相对均方误差的样本平均为 0.0348，中位数为 0.0115，分布明显偏斜。按 $v_R^{\rm agg}=v_R[\bar c+(1-\bar c)/J]$ 反推，州间环境的平均相关约为 $\bar c\approx0.34$。建议报告这一反推值及其 bootstrap 区间，说明样本量的限制，并可增加季节或滚动窗口以扩充样本。由于 $a$ 含有非环境成分，$\bar c$ 只能作粗略指示。依据见第~\ref{sec:empirical} 节。

\item 给覆盖率骤降补定量上界。均匀二项漏报下，增长率的离散程度至多放大 $\sqrt{\rho_1/\rho_2}=\sqrt{0.911/0.546}=1.29$ 倍，第 613 行的实测放大为 $1.823/0.624=2.92$ 倍。建议把这一对比写入正文，作为“结构性非随机扰动”判断的定量依据，与命题 3 的联系改为定性表述。依据见第~\ref{sec:empirical} 节。

\item 调整工具变量的定位。第 515 行的两阶段最小二乘给出 $b_{\rm iv}=2.5305$，高于 OLS 的 1.3525，这与回归元经典测量误差造成衰减的预期一致 \citep{fuller1987measurement,carroll2006measurement}。第一阶段 $F$ 统计量（未聚类 1147.1，按州聚类 312.9）只说明工具相关性充分；滞后 3--5 周的规模与流行阶段相关，排除限制无法检验，而随机分母对 $b$ 的放大发生在被解释变量中，工具变量无法消除。建议保持“替代测度下的稳健性”这一定位，在正文中讨论排除限制，并在中期报告与答辩口径中去掉“排除内生性”。依据见第~\ref{sec:empirical} 节。
\end{enumerate}

\subsubsection{写作与图表}

写作与图表方面的修改集中列于表~\ref{tab:ad-writing}。其中第 283 行的两阶段抽样类比建议补成一张符号对应表，形式可参照表~\ref{tab:ad-cochran}，并注明出自 Cochran 原书第 10 章 \citep{cochran1977sampling}。在无限总体下，Cochran 原式中单元间方差的位置对应 $v_R$；有限总体时该位置为 $S_1^2-S_2^2/M$，两种写法一致。

\begin{table}[!tbp]
\centering
\small
\caption{写作与图表的修改要点（行号指论文源码）}
\label{tab:ad-writing}
\begin{tabular}{L{2.6cm}L{5.8cm}L{6.4cm}}
\toprule
位置 & 问题 & 改法 \\
\midrule
第 756、761 行 & 附录“命题 5/6 的数学证明”两节标题与正文编号互换；超链接正确，标题文字错误 & 交换两节标题文字，答辩材料统一按正文编号 \\
第 251、312 行 & $A$ 同时指“目标 A”与方差系数 & 估计目标改称“目标 I（瞬时）”与“目标 II（跨期）”，或把方差系数改记为 $\kappa$ \\
第 254、330--332、418 行 & $r$ 同时指 AR(1) 系数与负二项形状参数 $nk$ & 形状参数统一写作 $nk$，或把 AR(1) 系数改记为 $\phi$ \\
第 176、291、409 行等 & 若干展示式缺少交叉引用标签，正文只能写“上式” & 被引用的式子补标签，不被引用的改为无编号环境 \\
第 187、196、399 行与中期报告 & 用词“拐点”不当；$V(m)$ 严格凸，没有拐点 & 统一改为“交叉点”或“平衡点”，并给出弹性 $-m_\times/(m+m_\times)$ \\
第 551、554 行 & 误差份额 54.8\%/45.2\% 与 54.6\%/45.1\% 的归一口径不同 & 第 551 行注明“按两项之和归一”，或统一采用第 554 行口径 \\
第 375 行 & $\max(1,I)$ 截断改变了小规模端的数据生成过程，触发频率未报告 & 报告各规模组的截断触发次数 \\
图 1(c)、图 4(c) & 图 1(c) 标题渲染为“m imes”；图 4(c) 的注释与柱条重叠 & 绘图脚本中的数学字符串改用原始字符串；调整注释位置 \\
图 1 至图 4 & 坐标轴与图例为英文，与中文正文不一致 & 改为中文标签 \\
\bottomrule
\end{tabular}
\end{table}

\begin{table}[!tbp]
\centering
\small
\caption{推论 1 与 Cochran 两阶段抽样的符号对应}
\label{tab:ad-cochran}
\begin{tabular}{L{6.2cm}L{6.2cm}}
\toprule
Cochran 两阶段抽样 & 本文推论 1 \\
\midrule
初级单元数 $n$ & 观测期数 $T$ \\
单元间方差 $S_b^2$ & 环境方差 $v_R$ \\
单元内方差 $S_w^2$ & 规模项系数 $A$ \\
每个初级单元的成本 $c_1'$ & 每期固定成本 $c_0$ \\
每个次级单元的成本 $c_2'$ & 每个报告病例的边际成本 $c_1$ \\
最优次级样本量 $m_{\rm opt}=\sqrt{(c_1'/c_2')(S_w^2/S_b^2)}$ & 最优单期规模 $m^*=\sqrt{(c_0/c_1)m_\times}$ \\
\bottomrule
\end{tabular}
\end{table}

\subsection{答辩与投稿问答准备}\label{sub:ad-qa}

下面 14 条追问按被问到的可能性与现有写法的薄弱程度选取。每条先概括追问内容，再给出建议的回答要点，最后注明依据位置。回答口径以第~\ref{sub:ad-revise} 小节的修改完成为前提；追问 1、2、8、9 最可能出现，应能用一两句话答出。

\begin{description}
\item[追问 1] 涉及 $m_\times=22.6$ 与理论区间 21--49 的“相符”能否算作对理论的验证。回答时应明确这一吻合不构成结构验证。实证 $a$ 是情景隐含值的 3.5--6.9 倍，$b$ 是 3.2--3.7 倍。$b$ 的放大约有一半来自用报告病例构造分母时引入的二项噪声，式~(\ref{eq:ad-denom}) 还能定量解释三种滞后权重下 $b$ 的排序；$a$ 的放大随回溯延迟单调增加，来自周增长比与世代再生数的口径差。两种机制不同，倍数接近属于巧合，所以 22.6 只说明量级一致，论文已改为这一口径。依据见第~\ref{sec:empirical} 节与第~\ref{sec:checks} 节。

\item[追问 2] 涉及 CRLB 的“有限样本收敛”指的是什么。回答时应先纠正措辞：估计量的方差在每个有限 $n$ 上精确等于 CRLB，不存在收敛过程。得分函数恰好是 $\mathcal I(R)(\widehat R-R)$，满足 Cauchy--Schwarz 不等式的取等条件，所以对任意 $n$ 都有 $\operatorname{Var}(\widehat R)=\mathcal I^{-1}$；图 1(b) 中 0.988--1.017 的偏离是蒙特卡洛误差。依据见第~\ref{sec:props} 节与附录引理~\ref{lem:ap-fisher}。

\item[追问 3] 涉及 $C/(\rho n)$ 为 UMVUE 的依据。回答要点为：固定 $nk$ 时 $C$ 的分布族是关于 $C$ 的单参数指数族，自然参数 $\ln\{\rho R/(k+\rho R)\}$ 关于 $R$ 严格单调，$C$ 是完备充分统计量；$C/(\rho n)$ 无偏，由 Lehmann--Scheffé 定理即得。UMVUE 性质以 $k$ 与 $\rho$ 已知为前提，完整证明在附录引理~\ref{lem:ap-fisher} 中给出。

\item[追问 4] 涉及周数据只有 26 周左右时 $T_{\rm eff}$ 近似的精度。回答要点为：$T_{\rm eff}$ 近似恒高估环境项，$T=26$、$r=0.8$ 时高估 20.5\%，所以正文已改用精确闭式，$T_{\rm eff}$ 只用于直觉说明。高估方差使区间估计偏保守；对设计而言，近似会略微偏向更多期数，$r=0.5$ 时近似最优单期规模为 42.1，精确值为 42.9。依据见第~\ref{sec:props} 节与图~\ref{fig:pb-ar1}。

\item[追问 5] 涉及称 $m_\times$ 为“拐点”时方差曲线在何处拐弯。回答要点为：$V(m)=v_R+A/m$ 的二阶导数 $2A/m^3>0$，曲线严格凸，没有拐点；$m_\times$ 是两类方差各占一半、弹性等于 $-1/2$ 的点。规模翻倍带来的方差降幅在小规模端接近 50\%，在 $m_\times$ 处为 25\%。论文已统一改称“交叉点”。依据见第~\ref{sec:props} 节与图~\ref{fig:pa-variance}。

\item[追问 6] 涉及离散检验的水平能否“严格低于 $\alpha$”。回答要点为：应写作“不超过 $\alpha$”，核验网格内实际水平最大达到 0.05000。非随机化检验只在它自己的实际水平下是一致最优；若需要恰好水平 $\alpha$，可以使用随机化检验，由 Karlin--Rubin 定理它是水平 $\alpha$ 的一致最优检验。依据见第~\ref{sec:props} 节。

\item[追问 7] 涉及图 2 中 80\% 功效样本量的稳定性。回答要点为：图中报告的是功效首次达到 0.8 的 $n$；临界值 $c_\alpha$ 随 $n$ 跳升，功效曲线呈锯齿，下一步就可能跌回 0.8 以下。此后恒不低于 0.8 的样本量多 1.4\%--3.0\%，论文已补报。依据见第~\ref{sec:props} 节与图~\ref{fig:pb-power}。

\item[追问 8] 涉及更新过程仿真称误差“不低于”单代基准，而图中按初始规模看低于基准。回答要点为：按初始规模排列确实低于基准（0.75--0.98 倍），因为增长期的实际规模中位数是初始值的 1.7--2.4 倍；按实现规模分箱后为 1.01--1.17 倍。分解式只保证偏差项非负，推不出一般下界，所以论文已改为经验陈述。依据见第~\ref{sec:props} 节。

\item[追问 9] 涉及 $\mathrm{CV}^2$ 随规模变化时 FGLS 采用 $1/\hat\mu^2$ 权重是否合理。回答要点为：均值设定正确时，任何只依赖 $D$ 的权重都给出一致估计，权重只影响效率；仿真验证了点估计无偏，州聚类 bootstrap 区间覆盖率为 0.944。更有效的权重是 $1/(\hat\mu^2\widehat{\mathrm{CV}}{}^2(D))$，已作为稳健性补充。依据见第~\ref{sec:empirical} 节与第~\ref{sec:checks} 节。

\item[追问 10] 涉及工具变量能否排除内生性。回答要点为：工具变量不能完全排除内生性。滞后规模与流行阶段相关，排除限制无法检验；工具变量主要处理回归元测量误差造成的衰减，所以 $b_{\rm iv}>b_{\rm ols}$ 符合预期。论文的定位是替代测度下规模项依然为正。依据见第~\ref{sec:empirical} 节。

\item[追问 11] 涉及提高报告率与扩大感染者基数两条路径的效果差距。回答要点为：同比例增加期望报告数时，两者的边际效果之比为式~(\ref{eq:ad-ratio}) 给出的 $A/\bar R$；住院监测中 $\rho$ 很低，这一比值只有 1.02--1.18，两条路径几乎等价，差异主要出现在高报告率、强超传播的场景。依据见第~\ref{sec:props} 节。

\item[追问 12] 涉及全国汇总为何能突破单州平台，以及这与 $v_R$ 不可约这一说法的关系。回答要点为：$v_R$ 只在同一环境内不可约；各州环境不完全相关时，汇总可削减其 $(1-\bar c)$ 部分。按全国相对均方误差 0.0348 反推 $\bar c\approx0.34$，但样本只有 28 个，只作粗略指示。依据见第~\ref{sec:empirical} 节。

\item[追问 13] 涉及 2024 年覆盖率骤降能否作为命题 3 的证据。回答要点为：覆盖率骤降只能作定性支持。均匀漏报下增长率的离散程度至多放大 1.29 倍，实测为 2.92 倍，说明主要来源是医院进出填报造成的结构性扰动，1.29 倍上界的推导在第~\ref{sec:empirical} 节。

\item[追问 14] 涉及推论 1 中 $m^*$ 区间的计算方法。回答要点为：$m^*=\sqrt{(c_0/c_1)m_\times}$，其中 $m_\times$ 取第 515 行流感情景的 21.3--49.2；这一结构与 Cochran 两阶段抽样同构，对应关系见表~\ref{tab:ad-cochran}。预算较小时 $T$ 的整数约束起作用，需要枚举 $T$。依据见第~\ref{sec:props} 节与图~\ref{fig:pb-budget}。
\end{description}

\subsection{投稿时的额外准备}

期刊审稿与学位答辩的关注点有所不同。除第~\ref{sub:ad-qa} 小节的问答外，投稿前建议额外做好以下准备。
\begin{itemize}
\item 审稿人很可能要求两向聚类与权重稳健性分析，建议在投稿前完成实证补充的第 1、2 项，并把结果放入正文或补充材料。
\item 周增长比与世代再生数的口径差应在方法节讲清，不宜留到局限性一节；截距 $a$ 的解读与三种权重下 $a$ 的单调变化可以一并给出。
\item 期刊版建议把式~(\ref{eq:ad-prop5}) 作为命题 5 的正文推论，并把随机分母机制写成独立的结果小节，使滞后权重敏感性同时具备理论预测与实证对照。
\item 精确 Wald 水平、稳定样本量与 $T_{\rm eff}$ 近似误差等数值细节，可放入补充材料，并附可复现的计算说明。
\end{itemize}

\subsection{后续研究方向}\label{sub:ad-future}

中期报告“存在问题及拟采取的措施”一栏列出了三个问题。方向 1 至方向 3 逐条给出与本讲义核验结果对接的做法，方向 4 至方向 6 为新增的研究方向，方向 7 处理研究计划本身。

\begin{description}
\item[方向 1] 处理空间相关性。中期报告称以“州级聚类稳健标准误控制截面自相关”，但按州聚类处理的是同州跨周的相关，同周跨州的截面相关并未处理。建议做两件事。其一是前述的两向聚类；其二是用 $v_R^{\rm agg}=v_R[\bar c+(1-\bar c)/J]$ 的框架估计州间环境相关 $\bar c$，可按地理邻接分层估计，与中期报告设想的空间滞后检验对接。这一做法比单纯加入空间权重矩阵更贴合论文的理论结构。

\item[方向 2] 检验亚型与平滑窗口的稳健性。中期报告拟论证“经验平衡点 $m_\times$ 在 14.6～24.4 例/周区间的稳健性”，其中 24.4 在论文中并不存在，三种滞后权重下的实际值为 12.64--45.96 例/周。建议把论证方式改为用式~(\ref{eq:ad-denom}) 预测各权重下的 $b$（包括中期报告提到的指数衰减权重），再检验预测与实证是否一致，敏感性由此从需要解释的问题变为被理论预测的结果。H1N1、H3N2 等优势亚型主导期的分段可同时检验 $a$ 是否随亚型变化。

\item[方向 3] 细化成本参数。推论 1 只依赖成本比 $c_0/c_1$，建议通过调研给出成本比的区间，配合整数 $T$ 枚举与 AR(1) 精确闭式，给出 $m^*$ 的敏感性表。$c_1$ 的口径需要统一：按报告病例计与按覆盖人群计相差一个因子 $\rho$，对应原文第 274 行的 $c_1=c_1'/\rho$。

\item[方向 4] 用双轨监测解耦截距 $a$。论文局限性一节已提到双轨监测。结合截距随回溯延迟增大的证据，可以设计周增长比与日度世代口径两套增长统计量，后者可用基于世代间隔的时变再生数方法估计 \citep{cori2013new}，从而直接估计口径差对 $a$ 的贡献。

\item[方向 5] 研究未知 $k$ 下的检验。命题 6 的精确水平以 $k$ 已知为前提。可用剖面似然或参数 bootstrap 校准临界值 $c_\alpha$，评估以 $\hat k$ 代入后实际水平与功效的变化。

\item[方向 6] 引入状态空间模型。式~(\ref{eq:ad-prop5}) 中的 $\mathrm{CV}^2(\Lambda_{t+1}\mid\mathcal F_t^C)$ 项可以借助状态空间模型的滤波方法估计，线性高斯近似下用 Kalman 滤波 \citep{hamilton1994time}，一般情形下用粒子滤波，从而把真实感染不可见的代价量化为可报告的数字。

\item[方向 7] 补齐研究计划的时间节点。中期报告“下一步研究工作计划”一栏为空白，需要补上各章定稿、投稿与答辩的时间节点。本章的建议条目可直接作为计划内容，建议的先后次序如下。先完成严谨性修改中的措辞与假设部分，它们不需要新的计算。再完成随机分母、两向聚类与权重对照等实证补充。最后处理写作与图表，并据此更新答辩材料。
\end{description}

\subsection{材料间遗留的不一致}

以下各条若已在优化后的中期报告中处理，答辩口径与优化后的中期报告保持一致即可；涉及论文仓库的，需要在仓库中修改。
\begin{itemize}
\item 论文仓库的说明文件（README）仍使用旧题“病例规模的推断收益边界：估计目标、监测杠杆与流感住院实证”，与现题目“疫情监测中的病例规模收益边界：估计目标、环境随机性与外部验证”不一致；说明文件称参考文献 38 条，参考文献库实际为 39 条。
\item 中期报告写“全美 50 州”，论文实证使用的是 50 个州及华盛顿特区（第 617 行），应补上华盛顿特区。
\item 中期报告所称的英文工作论文在仓库中找不到英文全文，答辩前应补齐全文或修改表述。
\item 中期报告称第 6、7 章“正准备展开正文起草”，与仓库中已有 7 章、85 页编译稿的状态不一致；“约 4.5 万字”无法从现有材料核实。
\item 中期报告用“微调”描述题目变更，与实际情况不符。病种由 COVID-19 改为流感，数据与题目也都已改变，建议如实说明变更内容与理由。
\item 论文类型填写为“基础研究”。论文同时包含理论推导、仿真与实证，可考虑改填“综合研究”。
\item 中期报告中“14.6～24.4”的区间与“排除内生性”的说法，按第~\ref{sub:ad-future} 小节方向 2 与实证补充第 6 项的口径更正。
\end{itemize}
